\documentclass[10pt,a4paper]{amsart}
\usepackage{amsmath,amssymb,amsthm,graphicx}
% Retained minimal preamble for the Stacks Project chapter more-algebra.tex (More on Algebra).
% The chapter file pinned-original-source.tex carries only \input{preamble} on its first line, so no
% preamble text lives inside the chapter itself.  The Stacks Project uses one single shared preamble for
% every chapter; the snapshot stages it identically as sources/stacks-chow/preamble.tex and
% sources/stacks-etale-cohomology/preamble.tex.  The block below is copied mechanically, line for line,
% from sources/stacks-chow/preamble.tex lines 14-18 (the Xy-pic setup) and 154-217 (the native theorem
% declarations together with the mathematical macros).  Nothing was added, removed or redefined; the
% wrapper class and the amsmath/amssymb/amsthm/graphicx base packages are supplied by the assembler.
% Chapter: more-algebra.tex -- More on Algebra.
% For commutative diagrams we use Xy-pic
\usepackage[all]{xy}

% We use 2cell for 2-commutative diagrams.
\xyoption{2cell}

% Theorem environments.
%
\theoremstyle{plain}
\newtheorem{theorem}[subsection]{Theorem}
\newtheorem{proposition}[subsection]{Proposition}
\newtheorem{lemma}[subsection]{Lemma}

\theoremstyle{definition}
\newtheorem{definition}[subsection]{Definition}
\newtheorem{example}[subsection]{Example}
\newtheorem{exercise}[subsection]{Exercise}
\newtheorem{situation}[subsection]{Situation}

\theoremstyle{remark}
\newtheorem{remark}[subsection]{Remark}
\newtheorem{remarks}[subsection]{Remarks}

\numberwithin{equation}{subsection}

% Macros
%
\def\lim{\mathop{\mathrm{lim}}\nolimits}
\def\colim{\mathop{\mathrm{colim}}\nolimits}
\def\Spec{\mathop{\mathrm{Spec}}}
\def\Hom{\mathop{\mathrm{Hom}}\nolimits}
\def\Ext{\mathop{\mathrm{Ext}}\nolimits}
\def\SheafHom{\mathop{\mathcal{H}\!\mathit{om}}\nolimits}
\def\SheafExt{\mathop{\mathcal{E}\!\mathit{xt}}\nolimits}
\def\Sch{\mathit{Sch}}
\def\Mor{\mathop{\mathrm{Mor}}\nolimits}
\def\Ob{\mathop{\mathrm{Ob}}\nolimits}
\def\Sh{\mathop{\mathit{Sh}}\nolimits}
\def\NL{\mathop{N\!L}\nolimits}
\def\CH{\mathop{\mathrm{CH}}\nolimits}
\def\proetale{{pro\text{-}\acute{e}tale}}
\def\etale{{\acute{e}tale}}
\def\QCoh{\mathit{QCoh}}
\def\Ker{\mathop{\mathrm{Ker}}}
\def\Im{\mathop{\mathrm{Im}}}
\def\Coker{\mathop{\mathrm{Coker}}}
\def\Coim{\mathop{\mathrm{Coim}}}

% Boxtimes
%
\DeclareMathSymbol{\boxtimes}{\mathbin}{AMSa}{"02}

%
% Macros for moduli stacks/spaces
%
\def\QCohstack{\mathcal{QC}\!\mathit{oh}}
\def\Cohstack{\mathcal{C}\!\mathit{oh}}
\def\Spacesstack{\mathcal{S}\!\mathit{paces}}
\def\Quotfunctor{\mathrm{Quot}}
\def\Hilbfunctor{\mathrm{Hilb}}
\def\Curvesstack{\mathcal{C}\!\mathit{urves}}
\def\Polarizedstack{\mathcal{P}\!\mathit{olarized}}
\def\Complexesstack{\mathcal{C}\!\mathit{omplexes}}
% \Pic is the operator that assigns to X its picard group, usage \Pic(X)
% \Picardstack_{X/B} denotes the Picard stack of X over B
% \Picardfunctor_{X/B} denotes the Picard functor of X over B
\def\Pic{\mathop{\mathrm{Pic}}\nolimits}
\def\Picardstack{\mathcal{P}\!\mathit{ic}}
\def\Picardfunctor{\mathrm{Pic}}
\def\Deformationcategory{\mathcal{D}\!\mathit{ef}}
\usepackage{cleveref}
\crefname{originallemma}{original lemma}{original lemmas}
\Crefname{originallemma}{Original lemma}{Original lemmas}
\makeatletter
\newlabel{derived-definition-cone}{{014E}{source}{}{sourcepointer-1}{}}
\newlabel{derived-definition-cone@cref}{{[originallemma][0][]014E}{[1][0][]source}{}{}{}}
\makeatother
\crefname{originallemma}{original lemma}{original lemmas}
\Crefname{originallemma}{Original lemma}{Original lemmas}
\makeatletter
\newlabel{derived-lemma-morphisms-from-projective-complex}{{064B}{source}{}{sourcepointer-2}{}}
\newlabel{derived-lemma-morphisms-from-projective-complex@cref}{{[originallemma][0][]064B}{[1][0][]source}{}{}{}}
\makeatother
\crefname{originallemma}{original lemma}{original lemmas}
\Crefname{originallemma}{Original lemma}{Original lemmas}
\makeatletter
\newlabel{derived-section-canonical-delta-functor}{{014Z}{source}{}{sourcepointer-3}{}}
\newlabel{derived-section-canonical-delta-functor@cref}{{[originallemma][0][]014Z}{[1][0][]source}{}{}{}}
\makeatother
\title{The Euler characteristic of a perfect complex is well defined on the derived category and additive on distinguished triangles}
\author{The Stacks Project Authors}
\date{stacks-project:more-algebra.tex:local-snapshot-20261012; GNU FDL 1.2 or later}
\begin{document}
\maketitle
\setcounter{subsection}{1}
\noindent\textit{Editorial note.} All delivered mathematical text is the author's own, copied from the pinned Stacks Project chapter file more-algebra.tex; no statement or proof step was rewritten, reordered or removed.\par\medskip
\noindent\textit{Editorial note.} The chapter file pinned-original-source.tex carries only \textbackslash{}input\{preamble\} on its first line, so no preamble text lives inside the chapter. The retained file sources/stacks-more-algebra/retained-preamble.tex is a mechanical, line-for-line copy of the project's single shared preamble as staged in the snapshot; nothing was added, removed or redefined and no package was added.\par\medskip
\noindent\textit{Editorial note.} The source URL is the source-level Stacks Project URL, as required when more than one card is produced from one source; the per-claim tag is carried by claim\_key.\par\medskip
\noindent\textit{Editorial note.} The author's displayed chain of short exact sequences contains a misprint in the second-to-last line, which reads '0 -> N\textasciicircum{}\{b-3\} -> M\textasciicircum{}\{b-2\} -> N\textasciicircum{}\{b-2\} -> 0' where the pattern of the neighbouring lines suggests N\textasciicircum{}\{b-2\}; the line is reproduced unchanged rather than silently corrected.\par\medskip
\noindent\textit{Allowed standard prerequisites (editorial).}\begin{itemize}
\item Perfect objects of D(R) and their representation by finite complexes of finite projective modules. used at proof lines 36508-36509. Supplies the representatives on which the formula for c is defined.
\item K\_0(R) and additivity of the class map on split short exact sequences of finite projective modules. used at proof lines 36512 and 36529. Is the target group and the relation that makes both the surjection step and the telescoping computation work.
\item Distinguished triangles of perfect complexes. condition (2) at statement lines 36499-36500. Fixes the additivity property that the constructed c must satisfy.
\end{itemize}\medskip
\ifdefined\hypertarget\hypertarget{sourcepointer-1}{}\fi
\noindent\textit{Source reference.} Original lemma 014E: focused-seeds/AGP/sources/stacks-snapshot-20261011/tags.txt, source lines 1469--1469. Cross-chapter navigation pointer into the Stacks Project; the standard fact it supplies is applied inside the counted proof span and the argument resting on it stays there.\par\medskip
\ifdefined\hypertarget\hypertarget{sourcepointer-2}{}\fi
\noindent\textit{Source reference.} Original lemma 064B: focused-seeds/AGP/sources/stacks-snapshot-20261011/tags.txt, source lines 7860--7860. Cross-chapter navigation pointer into the Stacks Project; the standard fact it supplies is applied inside the counted proof span and the argument resting on it stays there.\par\medskip
\ifdefined\hypertarget\hypertarget{sourcepointer-3}{}\fi
\noindent\textit{Source reference.} Original lemma 014Z: focused-seeds/AGP/sources/stacks-snapshot-20261011/tags.txt, source lines 1490--1490. Cross-chapter navigation pointer into the Stacks Project; the standard fact it supplies is applied inside the counted proof span and the argument resting on it stays there.\par\medskip
\begin{lemma}
\label{lemma-perfect-to-K-group}
Let $R$ be a ring. There is a map
$$
c : \text{perfect complexes over }R \longrightarrow K_0(R)
$$
with the following properties
\begin{enumerate}
\item $c(K[n]) = (-1)^nc(K)$ for a perfect complex $K$,
\item if $K \to L \to M \to K[1]$ is a distinguished triangle of
perfect complexes, then $c(L) = c(K) + c(M)$,
\item if $K$ is represented by a finite complex $M^\bullet$
consisting of finite projective modules, then
$c(K) = \sum (-1)^i[M_i]$.
\end{enumerate}
\end{lemma}

\begin{proof}
Let $K$ be a perfect object of $D(R)$. By definition we can represent
$K$ by a finite complex $M^\bullet$ of finite projective $R$-modules.
We define $c$ by setting
$$
c(K) = \sum (-1)^n[M^n]
$$
in $K_0(R)$. Of course we have to show that this is well defined,
but once it is well defined, then (1) and (3) are immediate.
For the moment we view the map $c$ as defined on complexes of
finite projective $R$-modules.

\medskip\noindent
Suppose that $L^\bullet \to M^\bullet$ is a surjective map
of finite complexes of finite projective $R$-modules.
Let $K^\bullet$ be the kernel. Then we obtain short exact
sequences of $R$-modules
$$
0 \to K^n \to L^n \to M^n \to 0
$$
which are split because $M^n$ is projective. Hence $K^\bullet$
is also a finite complex of finite projective $R$-modules and
$c(L^\bullet) = c(K^\bullet) + c(M^\bullet)$ in $K_0(R)$.

\medskip\noindent
Suppose given finite complex $M^\bullet$ of finite projective $R$-modules
which is acyclic. Say $M^n = 0$ for $n \not \in [a, b]$. Then we
can break $M^\bullet$ into short exact sequences
$$
\begin{matrix}
0 \to M^a \to M^{a + 1} \to N^{a + 1} \to 0, \\
0 \to N^{a + 1} \to M^{a + 2} \to N^{a + 3} \to 0, \\
\ldots \\
0 \to N^{b - 3} \to M^{b - 2} \to N^{b - 2} \to 0, \\
0 \to N^{b - 2} \to M^{b - 1} \to M^b \to 0
\end{matrix}
$$
Arguing by descending induction we see that $N^{b - 2}, \ldots, N^{a + 1}$
are finite projective $R$-modules, the sequences are split exact, and
$$
c(M^\bullet) = \sum (-1)[M^n] = \sum (-1)^n([N^{n - 1}] + [N^n]) = 0
$$
Thus our construction gives zero on acyclic complexes.

\medskip\noindent
It follows formally from the results of the preceding two
paragraphs that $c$ is well defined and satisfies (2). Namely,
suppose the finite complexes $M^\bullet$ and $L^\bullet$ of
finite projective $R$-modules represent the same object
of $D(R)$. Then we can represent the isomorphism by a map
$f : M^\bullet \to L^\bullet$ of complexes, see
Derived Categories, Lemma \ref{derived-lemma-morphisms-from-projective-complex}.
We obtain a short exact sequence of complexes
$$
0 \to L^\bullet \to C(f)^\bullet \to K^\bullet[1] \to 0
$$
see Derived Categories, Definition \ref{derived-definition-cone}.
Since $f$ is a quasi-isomorphism, the cone $C(f)^\bullet$ is
acyclic (this follows for example from the discussion in
Derived Categories, Section \ref{derived-section-canonical-delta-functor}).
Hence
$$
0 = c(C(f)^\bullet) = c(L^\bullet) + c(K^\bullet[1]) =
c(L^\bullet) - c(K^\bullet)
$$
as desired. We omit the proof of (2) which is similar.
\end{proof}

\end{document}
